\Needspace{8\baselineskip}
\Article{285}
Zu den Fragen, die unser drittes Problem bilden (auf das die beiden ersten
oben zurückgeführt wurden), können wir hier nur wenig sagen: Bei zwei
gegebenen ternären Formen derselben Determinante ist zu entscheiden,
ob sie äquivalent sind, und im ersten Fall sind alle Transformationen
der einen in die andere zu finden. Eine vollständige Lösung, wie wir sie
für die entsprechenden Probleme bei binären Formen gegeben haben, stößt
hier auf noch größere Schwierigkeiten.
\SourceLeaf{337}
Wir beschränken unsere Untersuchung daher auf einige besondere Fälle,
derentwegen dieser Exkurs hauptsächlich unternommen wurde.

I\quad Für die Determinante $+1$ haben wir oben gezeigt, dass alle
ternären Formen in zwei Klassen zerfallen, von denen die eine alle
indefiniten, die andere alle definiten (negativen) Formen enthält. Daraus
folgt sofort, dass zwei beliebige ternäre Formen mit der Determinante
$1$ äquivalent sind, wenn entweder beide definit oder beide indefinit
sind; ist dagegen die eine definit, die andere indefinit, so sind sie
nicht äquivalent (der letzte Teil gilt offenbar allgemein für Formen
jeder Determinante). --- Ebenso sind zwei beliebige Formen mit der
Determinante $-1$ sicher äquivalent, wenn entweder beide definit oder
beide indefinit sind. --- Zwei definite Formen mit der Determinante
$2$ sind stets äquivalent. Zwei indefinite sind nicht äquivalent,
wenn in der einen die ersten drei Koeffizienten alle gerade, in der
anderen aber nicht alle gerade sind; in den übrigen Fällen (wenn in
beiden die ersten drei Koeffizienten alle gerade sind oder in keiner)
sind sie äquivalent. --- Auf diese Weise könnten wir noch viel mehr
besondere Sätze angeben, wenn oben (Art.\ 277) mehr Beispiele ausgearbeitet
worden wären.

II.\ Bezeichnen $f$, $f'$ äquivalente ternäre Formen, so lässt sich
in allen diesen Fällen auch eine Transformation der einen in die andere
finden. Denn in allen Fällen haben wir oben in jeder Klasse ternärer
Formen eine recht kleine Anzahl von Formen angegeben, auf deren eine
sich jede Form derselben Klasse nach einheitlichen Methoden zurückführen
lässt; dort haben wir auch gezeigt, wie sich alle diese auf eine einzige
zurückführen lassen. Sei $F$ diese Form in der Klasse, in der $f$,
$f'$ liegen. Dann kann man nach den oben gegebenen Regeln Transformationen
der Formen $f$, $f'$ in $F$ und auch der Form $F$ in
$f$, $f'$ finden. Daraus lassen sich nach Art.\ 270 Transformationen
der Form $f$ in $f'$ und der Form $f'$ in $f$ ableiten.

III.\ Es bliebe also nur zu zeigen, wie man aus einer Transformation
der ternären Form $f$ in eine andere $f'$ alle möglichen
Transformationen ableiten kann. Dieses Problem hängt von einem einfacheren
ab, nämlich alle Transformationen der ternären Form $f$ in sich selbst
zu finden. Geht nämlich $f$ durch mehrere Substitutionen $(\tau)$,
$(\tau')$, $(\tau'')$ usw.\ in sich selbst und durch die Substitution $(t)$
in $f'$ über, so erhält man offenbar durch Verbindung der Transformation
$(t)$ mit $(\tau)$, $(\tau')$, $(\tau'')$ usw.\ nach Art.\ 270 Transformationen,
durch die $f$ jeweils in $f'$ übergeht. Ferner beweist man durch
Rechnung leicht, dass jede Transformation der Form $f$ in $f'$
auf diese Weise aus der Verbindung der gegebenen Transformation $(t)$
der Form $f$ in $f'$ mit irgendeiner (und zwar
\SourceLeaf{338}
\textit{genau einer}) Transformation der Form $f$ in sich selbst
abgeleitet werden kann. Verbindet man also die gegebene Transformation
der Form $f$ in $f'$ mit \textit{allen} Transformationen der Form
$f$ in sich selbst, so erhält man \textit{alle} Transformationen der
Form $f$ in $f'$, und zwar jede genau einmal.

Die Untersuchung aller Transformationen der Form $f$ in sich selbst
beschränken wir hier auf den Fall, dass $f$ eine definite Form ist,
deren Koeffizienten mit den Nummern $4,5,6$ alle $=0$ sind\HistoricalNote{*)}{Die übrigen Fälle, in denen $f$ eine definite Form ist, lassen sich auf diesen zurückführen. Ist dagegen $f$ eine indefinite Form, so muss man eine ganz andere Methode anwenden, und die Anzahl der Transformationen ist unendlich.}.
Sei also $f=\left(\begin{smallmatrix}a,\,a',\,a''\\0,\,0,\,0\end{smallmatrix}\right)$, und bezeichnen wir alle Substitutionen, durch die
$f$ in sich selbst übergeht, allgemein durch
\[
\begin{array}{ccc}
\alpha,&\beta,&\gamma\\
\alpha',&\beta',&\gamma'\\
\alpha'',&\beta'',&\gamma''
\end{array}
\]
sodass die Gleichungen
\[
\begin{aligned}
a\alpha\alpha+a'\alpha'\alpha'+a''\alpha''\alpha''&=a&&\ldots\quad(\Omega)\\
a\beta\beta+a'\beta'\beta'+a''\beta''\beta''&=a'\\
a\gamma\gamma+a'\gamma'\gamma'+a''\gamma''\gamma''&=a''\\
a\alpha\beta+a'\alpha'\beta'+a''\alpha''\beta''&=0\\
a\alpha\gamma+a'\alpha'\gamma'+a''\alpha''\gamma''&=0\\
a\beta\gamma+a'\beta'\gamma'+a''\beta''\gamma''&=0
\end{aligned}
\]
erfüllt sein müssen. Nun sind drei Fälle zu unterscheiden:

I.\ Sind $a$, $a'$, $a''$ (die dasselbe Vorzeichen haben)
alle verschieden, so nehmen wir $a<a'$, $a'<a''$ an (bei einer anderen
Größenfolge erhält man dieselben Schlüsse auf genau dieselbe Weise).
Dann verlangt die erste Gleichung in $(\Omega)$ offenbar $\alpha'=\alpha''=0$, also
$\alpha=\pm1$. Aus den Gleichungen $4,5$ folgt dann $\beta=0$, $\gamma=0$.
Ebenso folgt aus Gleichung $2$, dass $\beta''=0$, also $\beta'=\pm1$ ist;
aus Gleichung $6$ folgt $\gamma'=0$, und aus $3$, dass $\gamma''=\pm1$
ist. Da die Vorzeichen unabhängig voneinander gewählt werden können,
gibt es insgesamt $8$ verschiedene Transformationen.

II.\ Sind von den Zahlen $a$, $a'$, $a''$ zwei gleich,
\textit{z.B.}\ $a'=a''$, die dritte aber verschieden, so nehmen wir an,

\textit{erstens}, dass $a<a'$ ist. Dann gilt wie im vorigen Fall
$\alpha'=0$, $\alpha''=0$, $\alpha=\pm1$, $\beta=0$, $\gamma=0$. Aus den Gleichungen $2,3,6$
folgt aber leicht, dass
\SourceLeaf{339}
entweder $\beta'=\pm1$, $\gamma'=0$, $\beta''=0$, $\gamma''=\pm1$ oder $\beta'=0$, $\gamma'=\pm1$,
$\beta''=\pm1$, $\gamma''=0$ sein muss.

Ist dagegen \textit{zweitens} $a>a'$, so erhält man dieselben Schlüsse
folgendermaßen: Aus den Gleichungen $2,3$ folgt notwendig $\beta=0$,
$\gamma=0$, und entweder $\beta'=\pm1$, $\gamma'=0$, $\beta''=0$, $\gamma''=\pm1$ oder $\beta'=0$,
$\gamma'=\pm1$, $\beta''=\pm1$, $\gamma''=0$. Unter beiden Voraussetzungen folgt aus
den Gleichungen $4,5$, dass $\alpha'=0$, $\alpha''=0$ ist, und aus $1$,
dass $\alpha=\pm1$ ist. Für beide Fälle gibt es also $16$ verschiedene
Transformationen. --- Die übrigen beiden Fälle, in denen entweder
$a=a''$ oder $a=a'$ ist, erledigt man genauso, indem man die Zeichen
$\alpha$, $\alpha'$, $\alpha''$ im ersten Fall jeweils mit $\beta$,
$\beta'$, $\beta''$, im zweiten mit $\gamma$, $\gamma'$, $\gamma''$ vertauscht.

III.\ Sind alle $a$, $a'$, $a''$ gleich, so verlangen die
Gleichungen $1,2,3$, dass von den drei Zahlen $\alpha$, $\alpha'$,
$\alpha''$ sowie von $\beta$, $\beta'$, $\beta''$ und auch von $\gamma$,
$\gamma'$, $\gamma''$ jeweils zwei $=0$, die dritte $=\pm1$ sind.
Aus den Gleichungen $4,5,6$ sieht man aber leicht, dass von den drei
Zahlen $\alpha$, $\beta$, $\gamma$ nur eine $=\pm1$ sein kann;
ebenso von $\alpha'$, $\beta'$, $\gamma'$ und von $\alpha''$, $\beta''$,
$\gamma''$. Es gibt daher nur sechs Kombinationen:
\[
\begin{array}{c|c|c|c|c|c|c|}
\alpha&\alpha&\alpha'&\alpha'&\alpha''&\alpha''&=\pm1\\
\beta'&\beta''&\beta&\beta''&\beta&\beta'&=\pm1\\
\gamma''&\gamma'&\gamma''&\gamma&\gamma'&\gamma&=\pm1
\end{array}
\quad\text{Die übrigen sechs Koeffizienten }=0
\]
Wegen der Wahl der Vorzeichen erhält man insgesamt $48$ Transformationen.
--- Dasselbe Schema umfasst auch die vorigen Fälle: Von den ersten
sechs Spalten darf man aber nur die erste nehmen, wenn $a$,
$a'$, $a''$ alle verschieden sind; die erste und zweite, wenn
$a'=a''$; die erste und dritte, wenn $a=a'$; die erste und sechste,
wenn $a=a''$ ist.

Daraus folgt: Geht die Form $f=axx+a'x'x'+a''x''x''$ in eine andere äquivalente Form
$f'$ über, und zwar durch die Substitution
\[
x=\delta y+\varepsilon y'+\zeta y'',\quad x'=\delta'y+\varepsilon'y'+\zeta'y'',\quad x''=\delta''y+\varepsilon''y'+\zeta''y''
\]
so sind alle Transformationen
der Form $f$ in $f'$ in folgendem Schema enthalten:
\[
\begin{array}{c|c|c|c|c|c|l}
x&x&x'&x'&x''&x''&=\pm(\delta y+\varepsilon y'+\zeta y'')\\
x'&x''&x&x''&x&x'&=\pm(\delta'y+\varepsilon'y'+\zeta'y'')\\
x''&x'&x''&x&x'&x&=\pm(\delta''y+\varepsilon''y'+\zeta''y'')
\end{array}
\]
Dabei müssen alle ersten sechs Spalten verwendet werden, wenn $a=a'=a''$;
\SourceLeaf{340}
die Spalten $1$ und $2$, wenn $a'$, $a''$ gleich und
$a$ verschieden sind; $1$ und $3$, wenn $a=a'$; $1$
und $6$, wenn $a=a''$; schließlich nur die erste Spalte, wenn
$a$, $a'$, $a''$ alle verschieden sind. Im ersten Fall ist
die Anzahl der Transformationen $48$, im zweiten, dritten und vierten
$16$, im fünften $8$.

\Needspace{16\baselineskip}
\begin{center}
$\begin{array}{ccc}*&\quad&*\\&*&\end{array}$
\end{center}

\begin{center}
\textit{EINIGE ANWENDUNGEN AUF DIE THEORIE DER BINÄREN FORMEN.}
\end{center}

\begin{center}
\textit{Eine Form finden, durch deren Verdopplung eine gegebene binäre Form des Hauptgeschlechts entsteht.}
\end{center}

Von dieser kurzen Darstellung der ersten Grundlagen der Theorie ternärer
Formen gehen wir zu einigen besonderen Anwendungen über, unter denen
die folgende den ersten Platz verdient.

\Needspace{8\baselineskip}
\Article{286}
\textsc{Aufgabe.} \textit{Gegeben ist eine binäre Form $F=(A,B,C)$ mit der
Determinante $D$ aus dem Hauptgeschlecht. Gesucht ist eine
binäre Form $f$, durch deren Verdopplung sie entsteht.}

\textit{Lösung.} I.\ Man suche eine eigentliche Darstellung der zu $F$
entgegengesetzten Form $F'=ATT-2BTU+CUU$ durch die ternäre Form $xx-2yz$, etwa
\[
x=\alpha T+\beta U,\quad y=\alpha'T+\beta'U,\quad z=\alpha''T+\beta''U
\]
Dass dies möglich ist, folgt leicht aus der vorigen Theorie der ternären
Formen. Da nämlich $F$ nach Voraussetzung zum Hauptgeschlecht gehört,
gibt es einen Wert des Ausdrucks $\surd(A,\allowbreak B,\allowbreak C)(\mathrm{mod.}\,\allowbreak D)$. Daraus kann man eine ternäre
Form $\varphi$ mit der Determinante $1$ finden, die $(A,-B,C)$ als Teil
enthält. Dass alle Koeffizienten dieser Form ganze Zahlen sind, sieht
man ohne Schwierigkeit. Ebenso leicht sieht man, dass $\varphi$ eine
indefinite Form ist (denn nach Voraussetzung ist $F$ sicher keine
negative Form); sie ist daher notwendig zur Form $xx-2yz$ äquivalent.
Man kann also eine Transformation der letzten in die erste angeben,
die eine eigentliche Darstellung der Form $F'$ durch $xx-2yz$ liefert.
--- Dann gilt also
\[
A=\alpha\alpha-2\alpha'\alpha'',\quad -B=\alpha\beta-\alpha'\beta''-\alpha''\beta',\quad C=\beta\beta-2\beta'\beta''
\]
\SourceLeaf{341}
Bezeichnet man ferner die Zahlen $\alpha\beta'-\alpha'\beta$, $\alpha'\beta''-\alpha''\beta'$, $\alpha''\beta-\alpha\beta''$ mit
$a,b,c$, so haben sie keinen gemeinsamen Teiler, und es gilt $D=bb-2ac$.

II.\ Mithilfe der letzten Bemerkung aus Art.\ 235 folgt daraus leicht,
dass $F$ durch die Substitution $2\beta',\allowbreak \beta,\allowbreak \beta,\allowbreak \beta'';\ 2\alpha',\allowbreak \alpha,\allowbreak \alpha,\allowbreak \alpha''$ in das Produkt der Form
$(2a,-b,c)$ mit sich selbst und durch die Substitution $\beta',\allowbreak \beta,\allowbreak \beta,\allowbreak 2\beta'';\ \alpha',\allowbreak \alpha,\allowbreak \alpha,\allowbreak 2\alpha''$ in das
Produkt der Form $(a,-b,2c)$ mit sich selbst übergeht. Nun ist der größte
gemeinsame Teiler der Zahlen $2a,2b,2c$ gleich $2$. Ist also $c$
ungerade, so haben $2a,2b,c$ keinen gemeinsamen Teiler; mit anderen Worten, $(2a,-b,c)$
ist eine eigentlich primitive Form. Ebenso ist, wenn $a$ ungerade
ist, $(a,-b,2c)$ eine eigentlich primitive Form. Im ersten Fall entsteht
$F$ durch Verdopplung der Form $(2a,-b,c)$, im zweiten durch Verdopplung
der Form $(a,-b,2c)$ (vgl.\ Schluss 4, Art.\ 235); einer dieser Fälle muss
aber sicher stets eintreten. Wären nämlich beide $a,c$ gerade,
so müsste $b$ ungerade sein. Nun bestätigt man leicht, dass
$\beta''a+\beta b+\beta'c=0$, $\alpha''a+\alpha b+\alpha'c=0$ gilt; daraus würde folgen, dass $\beta b$, $\alpha b$,
also auch $\alpha$ und $\beta$ gerade sind. Dann wären aber $A$
und $C$ gerade, entgegen der Voraussetzung, dass $F$ eine
Form des Hauptgeschlechts und daher von eigentlich primitiver Ordnung
ist. --- Es kann übrigens auch sein, dass sowohl $a$ als auch
$c$ ungerade sind. In diesem Fall hat man sofort zwei Formen,
durch deren Verdopplung $F$ entsteht.

\textit{Beispiel.} Gegeben sei die Form $F=(5.\,2,31)$ mit der Determinante
$-151$. Als Wert des Ausdrucks $\surd(5,2,31)$ findet man hier $(55,22)$;
daraus folgt die ternäre Form $\varphi=\left(\begin{smallmatrix}5,&31,&4\\11,&0,&-2\end{smallmatrix}\right)$. Nach den Regeln aus Art.\ 272
findet man dazu die äquivalente Form $\left(\begin{smallmatrix}1,&1.&-1\\0,&0,&0\end{smallmatrix}\right)$, die in $\varphi$ durch die
Substitution $\left\{\begin{smallmatrix}2,&2,&1\\1,&-6,&-2\\0,&3,&1\end{smallmatrix}\right\}$ übergeht. Mithilfe der Transformationen aus Art.\ 277 findet
man dann, dass $\left(\begin{smallmatrix}1,&0,&0\\-1,&0,&0\end{smallmatrix}\right)$ in $\varphi$ durch die Substitution $\left\{\begin{smallmatrix}3.&-7,&-2\\2,&-1,&0\\1,&-9,&-3\end{smallmatrix}\right\}$
übergeht. Es ist also $a=11$, $b=-17$, $c=20$. Da $a$ ungerade
ist, entsteht $F$ durch Verdopplung der Form $(11,17,40)$ und geht
durch die Substitution $-1,-7,-7,-18;\ 2,3,3,2$ in das Produkt dieser Form mit sich
selbst über.

